\documentclass{amsart}
\usepackage{amsmath, amsthm, amssymb}
\newtheorem{theorem}{Question}
\title{FIM 590 - Homework 1}
\author{Andrew Lys}
\begin{document}
\maketitle
\begin{theorem}
    Write an HJB equation for 
    \[
    \inf_{\pi}\mathbb{E}\left[
        \frac{\gamma}{2} \int_0^T \pi_s^2 + X_T^2
    \right]
    \]
    where 
    \[
    d X_t = - \pi_t dt + \sigma d W_t
    \]
    with $X_0 = x \neq 0$ and $\gamma > 0$

    Solve the HJB using the ansatz 
    \[
    H(t,x) = a(t) x^2 + b(t)
    \]
\end{theorem}
\begin{proof}[Solution]
   We have 
   \[
   H(t, x) = \sup_\pi \mathbb{E}^x \left[
   -\frac{\gamma}{2} \int_t^T \pi_s^2 - X_T^2
   \right]
   \] 
   We have that the infinitesimal generator of $X$ is 
   \[
   \mathcal{L}^\pi = \frac{1}{2}\sigma^2 \partial^2_x -\pi_t \partial_x
   \]
   Then we have: 
   \[
   \begin{cases}
    \frac{\partial H}{\partial t} + \sup_{\pi} \left(
        \mathcal{L}^\pi H - \frac{\gamma}{2} \pi_t^2  
    \right) = 0\\
    H(T, x) = -x^2
   \end{cases}
   \]
   \begin{align*}
    \frac{\partial H}{\partial t} + \sup_{\pi} \left(
        \frac12 \sigma^2 \partial^2_x H - \pi_t \partial_x H - \frac{\gamma}{2} \pi_t^2  
    \right) = 0
   \end{align*}
   The supremum is attained at 
   \[
   \pi_t = -\frac{1}{\gamma} \partial_x H
   \]
   so we have
   \begin{align*}
    \frac{\partial H}{\partial t} + \frac{1}{2}\sigma^2 \partial^2_x H + \frac{1}{\gamma} (\partial_x H)^2 - \frac{1}{2\gamma} (\partial_x H)^2 = 0\\
    \frac{\partial H}{\partial t} + \frac{1}{2}\sigma^2 \partial^2_x H + \frac{1}{2\gamma} (\partial_x H)^2 = 0
   \end{align*}
   Plugging in the ansatz, we have: 
   \begin{align*}
   a'(t) x^2 + b'(t) + \sigma^2 a(t) + \frac{2}{\gamma}a(t)^2 x^2 = 0
   \end{align*}
   which gives us the following ODEs: 
   \[
   \begin{cases}
    a'(t) + \frac{2}{\gamma} a(t)^2 = 0\\
    b'(t) + \sigma^2 a(t) = 0\\
    a(T) = -1, b(T) = 0
   \end{cases}
   \]
   We can solve the first ODE to get: 
   \[
   a(t) = - \frac{\gamma}{2(T - t) + \gamma}
   \]
   and then solve the second ODE to get: 
   \[
   b(t) = -\frac{\sigma^2 \gamma}{2} \log{\left(
    \frac{2(T - t)}{\gamma} + 1
   \right)}
   \]
   and we have that the optimal control is: 
   \[
   \pi_t = \frac{1}{T - t + \gamma / 2} X_t
   \]
   to summarize: 
   \[
   \begin{cases}
    H(t, x) = - \frac{\gamma}{2(T - t) + \gamma} x^2 -\frac{\sigma^2 \gamma}{2} \log{\left(
    \frac{2(T - t)}{\gamma} + 1
   \right)}\\
   \pi_t = \frac{1}{T - t + \gamma / 2} X_t
   \end{cases}
   \]
\end{proof}
\begin{theorem}
    Consider a risky asset price:
    \[
    \frac{d S_t}{S_t} = \mu dt + \sigma dW_t
    \]
    for a portfolio with self-financing condition:
    \[
    \pi X_t^\pi = X_t^\pi \pi_t \frac{d S_t}{S_t} + r(1 - \pi_t) X_t^\pi dt - c_t dt
    \]
    $c_t \ge 0$. 

    Find the optimal consumption and the optimal value for 
    \[
    \sup_\pi \mathbb{E} \int_0^\infty e^{-\rho t} \log{(c_t)} dt
    \]
    where $\rho > 0$
\end{theorem}
\begin{theorem}
    Prove that Brownian motion has infinite total variation almost surely.
\end{theorem}
\begin{proof}[Solution]
    We prove a stronger statement: 
    Let $M$ be a continuous local martingale. 
    If $M$ is almost surely non-constant, then $M$ has infinite total variation almost surely.

    Let $V_M(t, \omega)$ be the total variation of $M$ on $[0,t]$.
    We assume for contradiction that $V_M(t) < \infty$ with positive probability.
    By stopping on the first time that $|M| > n$, we can assume that $M$ is bounded. 
    We then define
    \[
    T_0^n := \inf \{t > 0 : |M_t| > 2^{-n}\}, T_k^n := \inf \{t > T_{k-1}^n : |M_t - M(T_{k-1}^n)| > 2^{-n}\}, 
    \]
    and $t_k^n := T_k^n \wedge t$. 
    Then, we have that 
    \[
    [M]_t^n := \sum_{k \ge 1}(M(t_k^n) - M(t_{k-1}^n))^2 \to [M]_t 
    \]
    uniformly a.s. This is true by Rogers and Williams IV.30.1. (I don't want to rewrite the proof here)
    However, we also have that 
    \[
    [M]_t^n \le 2^{-n} \sum_{k \ge 1}|M(t_k^n) - M(t_{k - 1}^n)| \le 2^{-n} V_M(t)
    \]
    So, since with positive probability, we have that $V_M(t) < \infty$, then with positive probability, we have $[M]_t = 0$ for $t > 0$.
    But, since 
    \[
    M^2 - [M]
    \]
    is a martingale, stopping this martingale at 
    \[
    S:= \inf\{t > 0 : [M]_t > 0\}
    \] 
    and optionally stopping, we have 
    \[
    \mathbb{E}[M^2_{S \land t}] = \mathbb{E}[[M]_{S \land t}] = 0
    \]
    so that $M_{S \land t} = 0$, and since $S > 0$ with positive probability,
    We have that $M$ is constant with positive probability, which is a contradiction.
    
    Since Brownian Motion is a continuous local martingale, and is almost surely non-constant, we have that Brownian motion has infinite total variation almost surely.
\end{proof}
\end{document}